% ═══════════════════════════════════════════════════════════════════════════════
% CONCLUSIONES - INSTITUTO ATENEO DE COLIMA
% Estudio Documental: Corrupción en México (2006-2012)
% ═══════════════════════════════════════════════════════════════════════════════

\nota{
  Este capítulo sintetiza las metainferencias obtenidas mediante análisis 
  documental DIAC (CUAL + cuan) y contrasta hallazgos con literatura académica 
  (Casar y Cejudo 2024; Merino 2024; Sandoval Ballesteros 2024).
}

\section{Discusión, criterios de rigor y limitaciones del estudio}

La integración concurrente entre la evidencia documental dominante (CUAL) y las series estadísticas secundarias anidadas (cuan) permitió construir metainferencias consistentes con la literatura especializada \citep{hernandez2014}. En primer término, se observó plena convergencia entre la trayectoria descendente de México en el IPC (de 34 puntos en 2012 a 27 en 2025), el repunte del costo real de la corrupción cotidiana en la ENCIG 2025 (\$17,707.4 millones de pesos) y la evidencia cualitativa de los tres casos principales: la persistencia del fenómeno no obedece a la carencia de leyes, sino a la captura de controles horizontales, la discrecionalidad contractual y la impunidad transexenal descritas por Casar y Cejudo (2024), Merino (2024) y Sandoval Ballesteros (2024). En segundo lugar, las metainferencias mostraron que el encuadre mediático (framing) opera como un factor distorsionador de la rendición de cuentas cuando privilegia consignas virales imprecisas por encima de daños auditados y violaciones masivas a derechos humanos.

Para garantizar la calidad científica del reporte bajo un diseño DIAC de predominio cualitativo, se aplicaron cuatro criterios de rigor metodológico establecidos por Hernández Sampieri et al. (2014, cap. 14):

\begin{enumerate}
  \item \textbf{Credibilidad:} Se fortaleció mediante la triangulación sistemática de tres tipos de fuentes independientes ---resoluciones judiciales (DOJ, cortes de Florida, SCJN, CIDH), informes de fiscalización y derechos humanos (ASF, CNDH, COFEPRIS) y literatura académica mexicana e internacional arbitrada---, logrando corroboración estructural y adecuación referencial en cada caso.
  
  \item \textbf{Dependencia:} Se aseguró mediante la trazabilidad explícita de cada expediente (números de auditoría, expedientes judiciales, fechas exactas de emisión y enlaces verificables), criterios transparentes de inclusión muestral y el uso de instrumentos estandarizados de registro y cruce comparativo (Apéndices A y B), estableciendo cadenas de evidencia susceptibles de auditoría externa.
  
  \item \textbf{Confirmabilidad:} Para minimizar sesgos ideológicos o partidistas, se incorporaron expresamente las posturas de descargo del expresidente Felipe Calderón \citep{infobae2024,elfinanciero2025}, se documentó la continuidad transexenal de diversos entramados (contratos de García Luna hasta 2018, Odebrecht, Oceanografía y cancelación de Tula en 2014) y se integraron casos de contraste correspondientes a gobiernos de distinto signo político (Veracruz y Dos Bocas).
  
  \item \textbf{Transferencia analítica:} Se precisó que los hallazgos no buscan generalizarse estadísticamente a toda la Administración Pública Federal 2006--2012, sino ofrecer patrones institucionales analíticamente transferibles ---concentración de mando sin contrapesos, megainfraestructura sin candados transexenales y subrogación desregulada--- útiles para analizar otros periodos y niveles de gobierno.
\end{enumerate}

Asimismo, se reconocieron tres limitaciones metodológicas centrales que delimitan el alcance del estudio: (1) la dependencia exclusiva de fuentes documentales secundarias y archivos públicos, sin acceso a expedientes reservados ni entrevistas directas con actores judiciales o administrativos; (2) la vigencia de procedimientos aún en curso o en etapa de ejecución al cierre de noviembre de 2026 ---tales como el cobro efectivo de las sentencias civiles en Florida contra García Luna y su red empresarial, y el curso de acciones penales tras la declaratoria de imprescriptibilidad de la SCJN en el caso Guardería ABC---; y (3) la imposibilidad de realizar inferencias estadísticas hacia el universo completo de contrataciones y dependencias del sexenio 2006--2012 debido al carácter no probabilístico e intencional de la muestra (n = 3).

\section{Síntesis conclusiva y respuestas a las preguntas de investigación}

En conjunto, los tres casos principales aportaron evidencia documental congruente con la hipótesis de trabajo, con un alcance analíticamente transferible y no estadísticamente generalizable a toda la Administración Pública Federal 2006--2012. A partir de las metainferencias obtenidas, se dio respuesta puntual a las cinco preguntas de investigación:

\begin{enumerate}
  \item \textbf{Efectos verificables sobre el erario, la legitimidad y los derechos humanos:} La evidencia teórica y empírica examinada demostró que la corrupción ---desde prácticas cotidianas hasta esquemas sistémicos y de captura estatal--- produce efectos verificables y acumulativos sobre el erario (desvíos, sobrecostos y pagos indebidos que operan como impuesto regresivo), sobre la legitimidad democrática (deterioro sostenido de la confianza institucional y de la satisfacción con servicios públicos) y sobre el ejercicio de los derechos humanos fundamentales (vida, salud, seguridad y acceso equitativo a bienes públicos) \citep{cidh2019,inegi2026,monsivais2020,rios2021,roseackerman1999}.
  
  \item \textbf{Nivel y evolución empírica en México (2012--2026):} Los indicadores oficiales acreditaron la persistencia y el agravamiento de la corrupción en México más allá de un solo periodo gubernamental: la calificación nacional en el IPC descendió de 34 puntos al término de 2012 a 27 puntos en el IPC 2025, mientras que la ENCIG 2025 registró una percepción de frecuencia de 84.1\,\%, una prevalencia de 15.6\,\% en trámites y un costo directo de 17,707.4 millones de pesos anuales \citep{inegi2026,transparencia2026}.
  
  \item \textbf{Mecanismos institucionales y efectos de los tres casos principales:} En los tres expedientes del sexenio 2006--2012 se identificaron mecanismos estructurales convergentes: en el caso Genaro García Luna, la concentración sin contrapesos en la SSP se asoció con la captura de funciones de seguridad y desvíos contractuales sancionados en tribunales estadounidenses entre 2024 y 2026; en la Refinería Bicentenario, la ausencia de preinversión vinculante y candados transexenales dejó más de 9,000 millones de pesos federales sin planta productiva y una deuda estatal de 1,500 millones liquidada por Hidalgo hasta septiembre de 2025; y en la Guardería ABC, la expansión de la subrogación desregulada (81.5\,\% entre 2000 y 2009) junto con conflictos de interés y omisiones de inspección derivó en la muerte de 49 menores, cuya búsqueda de justicia alcanzó en junio de 2026 la declaratoria de imprescriptibilidad penal de la SCJN.
  
  \item \textbf{Brecha entre daño verificable y encuadre mediático:} El contraste con los episodios contemporáneos de Veracruz y Dos Bocas evidenció que la agenda mediática (framing) con frecuencia desplaza el escrutinio sobre daños masivos acreditados ---como 12 hospitales abandonados, 16.8 toneladas de medicamentos caducos o sobrecostos de 2.6 veces lo presupuestado--- hacia consignas virales imprecisas o polarizadas, debilitando la exigencia ciudadana de rendición de cuentas basada en pruebas.
  
  \item \textbf{Controles institucionales para reducir la brecha de implementación:} Con base en la literatura institucional y organizacional analizada \citep{arellanogault2012,casar2024,lopezayllon2009,merino2024,ocde2017,peschard2017}, contener la corrupción exige cerrar la brecha entre la ley y la sanción efectiva mediante cinco medidas concretas: garantizar la autonomía real de los órganos fiscalizadores y de procuración de justicia; reducir la discrecionalidad mediante la digitalización integral de trámites (apenas en 21.7\,\% en 2025); establecer candados técnicos vinculantes y transparencia abierta en megaproyectos transexenales; restituir la supervisión pública independiente y libre de conflictos de interés sobre servicios concesionados; y fortalecer la evaluación basada en indicadores verificables.
\end{enumerate}

En última instancia, observar la corrupción con la cabeza fría obliga a mirar más allá de las cifras presupuestales y los discursos coyunturales. Así como ningún proyecto duradero prospera cuando se improvisa y se abandonan los planes rigurosos, ningún país logra superar la impunidad si tolera que las reglas se doblen según conveniencias de grupo. Detrás de cada expediente documentado en estas páginas permaneció una lección difícil pero insoslayable: cuando se debilita la integridad de las instituciones, lo que realmente queda expuesto son la vida, la salud, la seguridad y el porvenir de las familias mexicanas. Reconocer con honestidad esos fallos constituye el primer paso indispensable para exigir controles que funcionen, sostener la persistencia en la rendición de cuentas y evitar que la historia vuelva a repetirse.